\newcommand{\duedate}{17 September 2026}
\documentclass{16384_doc}
\newcommand{\assignmentname}{Assignment 1: Forward Kinematics}

\begin{document}
\maketitle

\tableofcontents

\section{Code Questions}

In these problems you will implement the forward kinematics and workspace
analysis you derived in the written section, and then watch your own code drive
a simulated robot arm.

\subsection*{Setup}

Copy the code handout to a location of your choice. Everything below is run
from the top level of that folder --- the one containing
\verb!local_autograder.py!.

The simulator, the visualizer and the collision checker all live in
\verb!xarm7_lib!, and the environment file there installs them:

\begin{verbatim}
    conda env create -f xarm7_lib/environment.yml
    conda activate 16384
\end{verbatim}

The three exercises live in their own directories:

\begin{itemize}
    \item \verb!ex1! --- forward kinematics of a planar RRR arm.
    \item \verb!ex2! --- workspace analysis of a PR arm and of an RRR arm.
    \item \verb!ex3! --- the demo videos you record in the hands-on question.
\end{itemize}

\subsection*{Testing your work}

\verb!local_autograder.py! checks your functions against a large set of stored
test cases. Run all of them, or one exercise at a time:

\begin{verbatim}
    python local_autograder.py               # everything
    python local_autograder.py -p ex1        # just ex1
    python local_autograder.py -p ex2        # just ex2
\end{verbatim}

A failing run lists the first few cases that did not match; pass
\verb!--verbose! to see more. The stored cases live in \verb!exN/inK.txt!
(inputs) and \verb!exN/outK.txt! (expected outputs).

\textbf{A word of warning.} The autograder on Gradescope does not use those
stored outputs. It takes the same inputs, perturbs them by a small random
amount that is different on every submission, and compares your result against
a reference implementation. Code that reproduces the published numbers without
actually computing the kinematics will score zero there. Write the general
solution and the two graders will agree.

\begin{questions}

    \titledquestion{Forward Kinematics of an RRR Arm}[30]

Complete \verb!ex1/forward_kinematics_RRR.py!. When implemented correctly,\\
\verb!forward_kinematics_RRR! computes the forward kinematics of a planar RRR
arm from the three joint angles \verb!theta1!, \verb!theta2!, \verb!theta3! and
the three link lengths \verb!l1!, \verb!l2!, \verb!l3!.

Each link lies
along the $x$-axis of its starting frame, and each joint angle is measured
counter-clockwise from the $x$-axis of the previous link. Because the arm is
planar, every transform here is a $3 \times 3$ homogeneous matrix.

Seven frames are defined along the arm, alternating between a joint rotation
and a translation down a link:

\begin{itemize}
    \item \verb!H_1_0! $= \H{1}{0}$: the first joint, rotating \frame{0} by $\theta_1$ to give \frame{1}.
    \item \verb!H_2_1! $= \H{2}{1}$: the first link, translating \frame{1} by $l_1$ to give \frame{2}.
    \item \verb!H_3_2! $= \H{3}{2}$: the second joint, rotating \frame{2} by $\theta_2$ to give \frame{3}.
    \item \verb!H_4_3! $= \H{4}{3}$: the second link, translating \frame{3} by $l_2$ to give \frame{4}.
    \item \verb!H_5_4! $= \H{5}{4}$: the third joint, rotating \frame{4} by $\theta_3$ to give \frame{5}.
    \item \verb!H_6_5! $= \H{6}{5}$: the third link, translating \frame{5} by $l_3$ to give \frame{6}.
\end{itemize}

\frame{0} is the base frame of the whole robot and \frame{6} is the end
effector. Those six matrices are To-Do 1.

For To-Do 2, compose them into \verb!H_6_0! $= \H{6}{0}$, the transform taking a
point expressed in the end effector frame \frame{6} into the base frame
\frame{0}:
\[
    \H{6}{0} = \H{1}{0}\,\H{2}{1}\,\H{3}{2}\,\H{4}{3}\,\H{5}{4}\,\H{6}{5}
\]

The two To-Dos are marked separately, so a correct \verb!H_6_0! built from
incorrect elementary transforms will not earn full marks --- and neither will
correct elementary transforms that are never composed.

To check your implementation:
\begin{itemize}
    \item Start with angles of $0$ and $\piover{2}$, where you can work the
        answer out by hand, and confirm that each transform carries points from
        one link's frame into the previous one as you expect.
    \item Run \verb!python local_autograder.py -p ex1!.
    \item See it move: \verb!python local_autograder.py -p vis1! opens a browser
        with three sliders, one per joint angle. The arm follows the sliders,
        and the skeleton your \verb!forward_kinematics_RRR! predicts is drawn on
        top of it. When your kinematics are right the two lie exactly on top of
        one another; when they are wrong, the gap between them is your error.
        The hands-on question below says more about this.
\end{itemize}

    \titledquestion{Workspace Analysis of a PR Arm}[20]

Move to the \verb!ex2! directory. Here you will compute the reachable workspace
of a PR (prismatic-rotary) arm by sweeping its joints over their full ranges.

Start by reading \verb!workspace_analysis_RR.py!. It is a complete, working
example for an RR arm: \verb!endEffector_RR! places the end effector for one
pair of joint angles, and \verb!workspace_analysis_RR! sweeps both joints over
\verb!nSamples! values each and collects the $\mathtt{nSamples}^2$ resulting
positions. Compare the shape it produces to your analytical result from problem
8 of the written section. Use it as the template for everything that follows.

Your task is to complete \verb!endEffector_PR! in
\verb!workspace_analysis_PR.py!. The sweep itself is already written for you, so
this question is only the end effector position.

\begin{itemize}
    \item \verb!joints! is $[d, \theta]$: the prismatic extension $d$ in metres
        and the revolute angle $\theta$ in radians.
    \item \verb!linkLengths! is $[l_0, l_1]$. Note the convention here:
        \emph{$l_0$ is the base link length of the prismatic joint}, so the
        prismatic joint reaches from $l_0$ out to $l_0 + d$, and $l_1$ is the
        length of the rotary link beyond it.
    \item The prismatic joint ranges from $d = 0$ to $d = \d[f]$, its fully
        extended length. That range is passed in as \verb!minJoint[0]! and
        \verb!maxJoint[0]!.
\end{itemize}


Check your implementation with \verb!python local_autograder.py -p ex2!.

\titledquestion{Workspace Analysis of an RRR Arm}[30]

Now complete \verb!workspace_analysis_RRR.py!. Unlike the PR case, both
functions in this file are yours to write.

\begin{itemize}
    \item \verb!endEffector_RRR(joints, linkLengths)! returns the $(x, y)$
        position of the end effector, where \verb!joints! is
        $[\theta_1, \theta_2, \theta_3]$ and \verb!linkLengths! is
        $[l_1, l_2, l_3]$. This is the translation part of the \H{6}{0} you
        built in the forward kinematics question, and you should get the same
        answer both ways.
    \item \verb!workspace_analysis_RRR(nSamples, minJoint, maxJoint, linkLengths)!
        sweeps all three joints. Each joint takes \verb!nSamples! evenly spaced
        values between its entry in \verb!minJoint! and its entry in
        \verb!maxJoint!, and every combination is visited, so the returned
        \verb!xs! and \verb!ys! each hold $\mathtt{nSamples}^3$ entries.
\end{itemize}

The order in which you sweep the joints is up to you --- the autograder
compares the set of points you return, not the order they arrive in --- but the
number of points must be exactly $\mathtt{nSamples}^3$.

Check your implementation with \verb!python local_autograder.py -p ex2!.

\end{questions}

\section{Hands-On Questions}

\begin{questions}

    \titledquestion{Visualizing Your Kinematics on the xArm7}[20]

For this part of the homework you will run your own code on a simulated
UFACTORY xArm7 and record what you see. This question is graded by hand from
the videos you submit.

\subsection*{Why a 7-DoF arm can be your planar RRR arm}

The xArm7 has seven revolute joints; the arm you have been deriving all
assignment has three, in a plane. The two are the same mechanism once four of
the seven joints are held still. Holding
\[
    \theta_2 = \piover{2}, \quad
    \theta_3 = \piover{2}, \quad
    \theta_5 = -\piover{2}, \quad
    \theta_6 = \piover{2}
\]
leaves joints 1, 4 and 7 with their rotation axes all parallel to the world
$z$-axis. Everything that can still move is then confined to a horizontal
plane, and what is left is exactly a planar RRR arm: joint 1 at the shoulder,
joint 4 at the elbow, joint 7 at the wrist. The visualizer measures $l_1$,
$l_2$ and $l_3$ off the robot model itself and hands those lengths to your
functions, so the numbers you are tested on are the real geometry of the arm in
front of you.

\subsection*{Recording the videos}

Each command below opens a browser tab with the arm in it. Leave that tab
visible and record your screen while it run. Any common format is fine (\verb!.mp4!, \verb!.mov!,
\verb!.webm!).

\begin{enumerate}
    \item \textbf{Forward kinematics.}

\begin{verbatim}
    python local_autograder.py -p vis1
\end{verbatim}

    Three sliders appear above the arm, one per joint angle, and the skeleton
    predicted by your \verb!forward_kinematics_RRR! is drawn over the arm as you
    drag them. Record yourself moving all three sliders through a good part of
    their range. Your video should show the predicted skeleton staying on the
    arm throughout. If it drifts off, your kinematics are wrong --- fix them
    before recording.

    \item \textbf{Workspace sweep.}

\begin{verbatim}
    python local_autograder.py -p vis2
\end{verbatim}

    The points your \verb!workspace_analysis_RRR! predicts are drawn as a faint
    cloud, and the arm then walks through the same joint grid, lighting up each
    point as it reaches it. Record enough of the sweep that the traced points
    are clearly filling in the predicted cloud.

    The joint limits for this sweep are chosen so that the arm never collides
    with itself, which is why the elbow angle does not use its full range.
\end{enumerate}

\subsection*{Deliverables}

Put both videos in the \verb!ex3! directory; \verb!create_submission.py! will
pick up everything in there. 

Run \verb!create_submission.py! from the top level of the handout. It collects
your three completed Python files and everything in \verb!ex3!, and writes a
single \verb!handin.zip!:

\begin{verbatim}
    python create_submission.py
\end{verbatim}


This script does not grade your submission; it only packages it. Use
\verb!local_autograder.py! to check your answers before you submit.

\end{questions}

\begin{submissionChecklist}
    \item Complete \verb!ex1/forward_kinematics_RRR.py!,
        \verb!ex2/workspace_analysis_PR.py! and
        \verb!ex2/workspace_analysis_RRR.py!.
    \item Confirm \verb!python local_autograder.py! passes.
    \item Record the two demo videos and put them in \verb!ex3!.
    \item Run \verb!python create_submission.py! to produce \verb!handin.zip!.
    \item Upload \verb!handin.zip! to the \textbf{HW1 Code} assignment on
        Gradescope.
\end{submissionChecklist}

\end{document}
